% NUST_presentation_template.tex -- beamer template for MS/PhD defence presentation at NUST
% (proposal defence, pre-defence, final defence).
% infolines outer theme,
% NUST red/grey colour scheme, IEEE-style numeric references, and one frame
% layout per section of the usual defence outline. Replace the placeholder
% text; keep or drop frames as the talk needs.
%
% Build:  latexmk -pdf NUST_presentation_template.tex   (needs biber)
% Bibliography: references.bib next to this file (replace with your own).

% ---------------------------------------------------------------- aspect ratio
% Change the number after aspectratio= and rebuild; nothing else needs editing.
%   43   4:3   (128 x 96 mm)  -- the classic near-square projector format (default here)
%   169  16:9  (160 x 90 mm)  -- widescreen; the format the layouts were first drawn for
%   1610 16:10, 149 14:9, 141 1.41:1, 54 5:4 (the squarest), 32 3:2
% Every width in the file is a fraction of \textwidth and every placeholder height a
% fraction of \textheight, so the frames rescale; at the narrower ratios check the
% four-column tables and the two-column result frames for crowding and drop to
% \scriptsize or \tiny where needed.
\documentclass[aspectratio=43, 11pt, xcolor=dvipsnames, compress]{beamer}

% ---------------------------------------------------------------- theme
\mode<presentation>
\usefonttheme{professionalfonts}
\useinnertheme{rectangles}
\useoutertheme{infolines}
\setbeamertemplate{navigation symbols}{}
\setbeamertemplate{caption}{\raggedright\insertcaption\par}   % no "Figure n:" on slides
\setbeamerfont{caption}{size=\scriptsize}
\setbeamertemplate{itemize items}[circle]
\setbeamertemplate{enumerate items}[default]
\setbeamertemplate{bibliography item}{\insertbiblabel}          % [n] labels in the reference list

% colours -- NUST red/grey scheme
\definecolor{nustred}{RGB}{160,30,35}
\definecolor{nustgrey}{RGB}{60,60,60}
\setbeamercolor{structure}{fg=nustred}
\setbeamercolor{title}{fg=White,bg=nustred}
\setbeamercolor{frametitle}{fg=nustgrey,bg=Black!8}
\setbeamercolor{framesubtitle}{fg=nustred}
\setbeamercolor{section in head/foot}{fg=White,bg=nustred!85}
\setbeamercolor{subsection in head/foot}{fg=nustred,bg=Black!12}
\setbeamercolor{title in head/foot}{fg=nustgrey,bg=Black!8}
\setbeamercolor{author in head/foot}{fg=White,bg=nustred}
\setbeamercolor{date in head/foot}{fg=nustgrey,bg=Black!15}
\setbeamercolor{block title}{fg=White,bg=nustred!85}
\setbeamercolor{block body}{bg=nustred!5}
\setbeamercolor{block title alerted}{fg=White,bg=nustgrey}
\setbeamercolor{block body alerted}{bg=Black!5}
\setbeamercolor{alerted text}{fg=nustred}

\setbeamerfont{frametitle}{size=\large,series=\bfseries}
\setbeamerfont{framesubtitle}{size=\small}
\setbeamerfont{footnote}{size=\tiny}

% ---------------------------------------------------------------- packages
\usepackage[utf8]{inputenc}
\usepackage[T1]{fontenc}
\usepackage{lmodern}
\usepackage[english]{babel}
\usepackage{amsmath, amssymb, amsfonts, mathtools}
\usepackage{braket}
\usepackage{graphicx}
\graphicspath{{./}{figures/}}
\usepackage{booktabs}
\usepackage{tabularx}
\usepackage{multirow}
\usepackage{adjustbox}   % max width=\linewidth shrinks a table that is wider than its column
\usepackage{tikz}
\usepackage{xcolor}
\usepackage{csquotes}

% ---------------------------------------------------------------- bibliography (optional)
% If references.bib exists next to this file it is used (biblatex + biber, IEEE numeric
% style). If it does not, biblatex is not loaded, every \cite prints nothing, and the
% References frame at the end is skipped, so the file builds with pdflatex alone and no
% biber run. To use a bibliography under another name, change the filename in both
% places below.
\newif\ifhasbib
\IfFileExists{./references.bib}{\hasbibtrue}{\hasbibfalse}
\ifhasbib
  \usepackage[backend=biber, style=ieee, citestyle=numeric-comp, doi=false, isbn=false, url=false, eprint=false, maxbibnames=3, minbibnames=1]{biblatex}
  \addbibresource{references.bib}
  \renewcommand*{\bibfont}{\scriptsize}
  \setlength{\bibitemsep}{0.25ex}
\else
  \renewcommand{\cite}[2][]{\unskip}   % swallow the citation and the space before it
\fi

% column types for the wide comparison tables
\newcolumntype{Y}{>{\raggedright\arraybackslash}X}
\newcolumntype{F}{>{\raggedright\arraybackslash\bfseries}p{0.11\textwidth}}

% +/- markers for the "how the research questions are answered" tables
\newcommand{\rqplus}{\alert{$+$}}
\newcommand{\rqminus}{\textcolor{nustgrey}{$-$}}

% grey placeholder box standing in for a figure:
%   \placeholderfig{<width as a fraction of the current \textwidth>}{<height as a fraction of \textheight>}{caption}
% Inside a \column, \textwidth is the column width, so {1}{0.5} fills the column and
% takes half the frame height whatever the aspect ratio is.
\newcommand{\placeholderfig}[3]{%
  \begin{tikzpicture}
    \draw[fill=Black!6, draw=Black!30] (0,0) rectangle (#1\textwidth,#2\textheight);
    \node[text=nustgrey, font=\scriptsize, align=center, text width=0.9*#1\textwidth] at (0.5*#1\textwidth,0.5*#2\textheight) {#3};
  \end{tikzpicture}}

% ---------------------------------------------------------------- title page
% Edit the four bracketed items. The title, author and institute come from
% \title / \author / \institute below.
\setbeamertemplate{title page}{
  \vfill
  \centering
  \usebeamertemplate{title}
  \vskip0.2cm
  \usebeamertemplate{titlegraphic}
  \vskip0.2cm
  \usebeamertemplate{author}
  \vskip0.1cm
  {\small [Pre-Defence / Final Defence]}\\[0.1cm]
  {\footnotesize [MS / PhD] [Name of degree]}\\[0.25cm]
  {\normalsize Supervisor: [Prof.] Dr.\ [Name of supervisor]}\\[0.25cm]
  {\usebeamerfont{institute}[Department], School of Natural Sciences\\
   National University of Sciences and Technology (NUST), Islamabad, Pakistan}\\[0.15cm]
  {\usebeamerfont{date}\insertdate}
  \vfill
}

\title[Short title]{Full Title of the Thesis,\\Possibly on Two Lines}
\author[Student name]{Student name}
\institute[SNS, NUST]{School of Natural Sciences, NUST}
\date{\today}
\titlegraphic{\includegraphics[height=0.15\textheight]{nust-logo.png}}

% ================================================================
\begin{document}

\begin{frame}[plain]
  \titlepage
\end{frame}

% Two-column outline. The ranges must cover the numbered sections below
% (ten as written); adjust them if sections are added or removed.
\begin{frame}{Outline}
  \small
  \begin{columns}[T]
    \column{0.48\textwidth}
      \tableofcontents[sections={1-5}]
    \column{0.48\textwidth}
      \tableofcontents[sections={6-10}]
  \end{columns}
\end{frame}

% ================================================================
\section{Introduction}

% Two-column background frame: the classical / prior picture on the left,
% the proposed one on the right, then one block stating why it matters.
\begin{frame}{Background: [from the established picture to the proposed one]}
  \small
  \begin{columns}[T]
    \column{0.5\textwidth}
      \textbf{[The established approach]}
      \begin{itemize}
        \item what it does, in one line; an equation if there is one: $x_{t+1} = f(x_t)$
        \item its characteristic behaviour
        \item where it is used
      \end{itemize}
    \column{0.5\textwidth}
      \textbf{[The proposed approach]}
      \begin{itemize}
        \item what changes, with a citation~\cite{example_article}
        \item the defining relation,
          \[ U = S\,(\mathbb{I}\otimes C) \]
        \item its characteristic behaviour, contrasted with the left column
      \end{itemize}
  \end{columns}
  \vskip0.3cm
  \begin{block}{Why this is a candidate [general tool / better method]}
    One or two sentences, each supported by a citation~\cite{example_book, example_inproceedings}, giving the reason the proposed approach deserves the study.
  \end{block}
\end{frame}

\begin{frame}{Motivation}
  \small
  \begin{itemize}
    \item \textbf{First motivating fact.} State it and say why it matters.
    \item \textbf{Second motivating fact.} The step from the first to the problem studied.
    \item \textbf{Third.} What changes when the proposed approach is adopted, with a citation~\cite{example_article}.
  \end{itemize}
  \begin{alertblock}{The practical interest}
    Where the results would be used in practice, and why a method with these properties is plausibly a common ingredient of those uses.
  \end{alertblock}
\end{frame}

% ================================================================
\section{Literature Review}

% Two-column review: one column per parent field.
\begin{frame}{Literature review (1): [first parent field]; [second parent field]}
  \small
  \begin{columns}[T]
    \column{0.5\textwidth}
      \textbf{[First field]}
      \begin{itemize}
        \item key result one~\cite{example_article}
        \item key result two~\cite{example_inproceedings}
        \item key result three
        \item foundational theorems~\cite{example_book}
      \end{itemize}
    \column{0.5\textwidth}
      \textbf{[Second field]}
      \begin{itemize}
        \item key result one~\cite{example_article}
        \item key result two; the known obstacle~\cite{example_inproceedings}
        \item benchmarking and the state of the evidence~\cite{example_book}
      \end{itemize}
  \end{columns}
\end{frame}

% The intersection of the two fields: the work closest to this thesis.
\begin{frame}{Literature review (2): [the intersection of the two fields]}
  \small
  A considerably smaller body of work than either parent field.
  \begin{itemize}
    \item \textbf{[Closest antecedent one.]} What it did~\cite{example_article} --- and in what respect it differs from this work.
    \item \textbf{[Closest antecedent two.]} What it did~\cite{example_inproceedings} --- and what it is \emph{not}.
    \item \textbf{[Design choices that matter and are seldom varied.]} Evidence that a choice usually fixed by convention changes the outcome~\cite{example_book}.
  \end{itemize}
\end{frame}

\subsection{Research Gaps Identified}
\begin{frame}{Research gaps identified}
  \small
  \begin{block}{Gap 1 --- [short name]}
    What the literature asserts or assumes, and what it has not actually done.
  \end{block}
  \begin{block}{Gap 2 --- [short name]}
    Which free choices are fixed by convention, with evidence that they matter~\cite{example_article}.
  \end{block}
  \begin{block}{Gap 3 --- [short name]}
    What comparison is missing, and under what protocol it would have to be made~\cite{example_inproceedings}.
  \end{block}
\end{frame}

% ================================================================
\section{Problem Statement}
\begin{frame}{Problem statement}
  \footnotesize
  \begin{itemize}
    \item The observation the problem rests on: \textbf{what is seldom varied, measured, or compared} in existing work.
    \item Why the obvious justification for the proposed approach is not by itself an answer.
  \end{itemize}
  \begin{alertblock}{The problem addressed}
    One sentence stating exactly what is instantiated, on what terms, across which cases ---
    \textbf{[case family A]} and
    \textbf{[case family B]} ---
    and what is measured for each.
  \end{alertblock}
  \vskip0.1cm
  \textbf{Framing:} the one-line framing of the thesis, and what remains the application.\\
  \textbf{Limits fixed at the outset:} what is and is not claimed; the simulation or experimental regime; how a negative outcome is treated.
\end{frame}

% ================================================================
\section{Research Questions}
\begin{frame}{Research questions}
  For every case the same three questions are posed:
  \vskip0.3cm
  \begin{enumerate}
    \item[\textbf{RQ1}] \textbf{[Formulation.]} What is asked first, stated precisely.
    \vskip0.2cm
    \item[\textbf{RQ2}] \textbf{[Cost.]} What is asked second, and in which units.
    \vskip0.2cm
    \item[\textbf{RQ3}] \textbf{[Performance.]} What is asked third, and against what baseline.
  \end{enumerate}
  \vskip0.4cm
  \begin{block}{Variables of the investigation (not its subject)}
    The parameters that are varied wherever a case makes the comparison meaningful, listed so the committee can see they are controls rather than the contribution.
  \end{block}
\end{frame}

% ================================================================
\section{Research Objectives}
\begin{frame}{Research objectives}
  \begin{enumerate}
    \item[\textbf{O1}] Produce a \textbf{[deliverable one]} on [cases].
    \item[\textbf{O2}] Exhibit \textbf{[deliverable two]}, simulate it, and cost it.
    \item[\textbf{O3}] Show that \textbf{[claim three]}, and evaluate it against [baselines].
    \item[\textbf{O4}] Reproduce \textbf{[known result]} and quantify [what it costs].
    \item[\textbf{O5}] Build \textbf{[the learning / application models]} and evaluate them under a fixed protocol against tuned baselines.
    \item[\textbf{O6}] Report negative results with the same precision as positive ones, and identify the \textbf{regime in which the approach is advantageous}.
  \end{enumerate}
\end{frame}

% ================================================================
\section{Methodology}

\subsection{Justification of Proposed Methodology}
\begin{frame}{Justification of the methodology}
  \small
  \begin{itemize}
    \item \textbf{Why [the central construction] everywhere?} The reason, with a citation~\cite{example_article}; where the alternative construction is still used and why. (Note any committee observation from the proposal defence that this answers.)
    \item \textbf{Why one primitive across [several] cases?} What that is the only way to test.
    \item \textbf{Why [simulation / this experimental setting]?} What the setting guarantees and what it leaves out.
    \item \textbf{Why [this reporting discipline]~\cite{example_inproceedings}?} Tuned baselines, splits fixed before model selection, results over seeds.
    \item \textbf{Why report cost in [these units]?} What the more common unit hides.
  \end{itemize}
\end{frame}

\subsection{Alternative Approaches}
\begin{frame}{Alternative approaches considered}
  \footnotesize
  \renewcommand{\arraystretch}{1.15}
  \begin{tabularx}{\textwidth}{@{}>{\raggedright\arraybackslash\bfseries}p{0.30\textwidth} >{\raggedright\arraybackslash}X@{}}
    \toprule
    Alternative & Why it was not the main route \\
    \midrule
    {}[Alternative one]~\cite{example_article} & Why it does not answer the research questions; what is built instead. \\
    {}[Alternative two] & Its cost or prerequisite; where it is nevertheless used. \\
    {}[Standard constructions from the literature] & Kept --- but as tuned \emph{baselines}, not as the object of study. \\
    {}[Runs on real hardware / field data] & Not needed to answer RQ1--RQ3; what would obscure the comparison. Future work. \\
    {}[Out-of-scope generalisation]~\cite{example_book} & Out of scope; the restriction that applies throughout. \\
    \bottomrule
  \end{tabularx}
\end{frame}

\subsection{Details}
% Two-column details frame: definitions on the left, the axes of variation
% on the right.
\begin{frame}{Details (1): [the central object as a design space]}
  \footnotesize
  \begin{columns}[T]
    \column{0.55\textwidth}
      \textbf{Setting.} The objects, their sizes, and the notation used for the rest of the talk.
      \vskip0.2cm
      \textbf{The family used throughout:}
      \[
        K(\delta,\theta) = \Phi(\delta)\,e^{-i\theta G}
      \]
      \begin{itemize}
        \item what the first parameter interpolates between
        \item what the second parameter interpolates between
        \item a property that makes it usable in the method (e.g.\ differentiable)
      \end{itemize}
    \column{0.45\textwidth}
      \textbf{Two axes of variation}
      \begin{enumerate}
        \item \textbf{[Axis one]}: what it depends on; its cost.
        \item \textbf{[Axis two]}: what it depends on; parameter count; the axis the trained models use.
      \end{enumerate}
      \vskip0.2cm
      \textbf{Characterised before any objective is optimised:} the quantities measured for every configuration.
  \end{columns}
\end{frame}

% Four-column construction / baseline / protocol table, one row per case.
\begin{frame}{Details (2): [first group of cases] --- construction, baseline, protocol}
  \footnotesize
  \renewcommand{\arraystretch}{1.2}
  \begin{tabularx}{\textwidth}{@{}F Y Y Y@{}}
    \toprule
    Case & Construction & Baseline & Protocol \\
    \midrule
    {}[Case A] & what is built & what it is compared with & sizes, seeds, tooling, repetitions \\
    {}[Case B] & what is built & what it is compared with & sizes, seeds, tooling, repetitions \\
    {}[Case C] & what is built & what it is compared with & sizes, seeds, tooling, repetitions \\
    {}[Case D] & what is built & what it is compared with & sizes, seeds, tooling, repetitions \\
    \bottomrule
  \end{tabularx}
\end{frame}

\begin{frame}{Details (3): [second group of cases] --- construction, baselines, protocol}
  \footnotesize
  \renewcommand{\arraystretch}{1.2}
  \begin{tabularx}{\textwidth}{@{}F Y Y Y@{}}
    \toprule
    Case & Construction & Baselines & Protocol \\
    \midrule
    {}[Case E] & what is built & the baselines, tuned~\cite{example_article} & dataset, split, loss, optimiser, budget, seeds \\
    {}[Case F] & what is built & the baselines, tuned & datasets~\cite{example_inproceedings}; cross-validation; splits \\
    {}[Case G] & what is built & the baselines, tuned~\cite{example_book} & datasets; folds \\
    \bottomrule
  \end{tabularx}
  \vskip0.2cm
  One sentence on the protocol rule that applies to every row (e.g.\ hyperparameters selected on the training fold alone).
\end{frame}

% ================================================================
\section{Implementation}
\begin{frame}{Implementation}
  \footnotesize
  \begin{columns}[T]
    \column{0.5\textwidth}
      \textbf{Tooling}
      \begin{itemize}
        \item libraries and versions
        \item any custom core, and how it was cross-checked
        \item validation against an independent reference
      \end{itemize}
      \textbf{Reproducibility}
      \begin{itemize}
        \item every stochastic element seeded and stated
        \item results written to a record; every table and figure generated from it by script
      \end{itemize}
    \column{0.5\textwidth}
      \textbf{Code organisation} (\texttt{[repository]})
      \begin{itemize}
        \item \texttt{[module one]} --- what it holds
        \item \texttt{[module two]} --- what it holds
        \item \texttt{[module three]} --- what it holds
      \end{itemize}
      \textbf{Limits}
      \begin{itemize}
        \item the size or regime that is tractable, and what sets it
        \item what is idealised away, and what that makes the reported numbers (e.g.\ a ceiling)
      \end{itemize}
  \end{columns}
\end{frame}

% Figure frame: two stacked figures with a short reading below.
% Replace \placeholderfig with \includegraphics[width=0.8\textwidth, height=0.36\textheight, keepaspectratio]{figure.pdf}.
\begin{frame}{Implementation: [the central construction as a diagram]}
  \centering
  \placeholderfig{0.8}{0.32}{figure 1 --- e.g.\ the construction as a circuit or block diagram}\\[0.1cm]
  \placeholderfig{0.8}{0.22}{figure 2 --- e.g.\ one component in detail}
  \begin{itemize}\small
    \item What one layer / unit of the construction consists of.
    \item The feature of the construction the results will turn on.
  \end{itemize}
\end{frame}

% ================================================================
\section{Results and Evaluation}

% Result frame, layout A: figure left, table + bullets right, one-line
% caveat at the bottom. Replace \placeholderfig with \includegraphics.
\begin{frame}{Result 1 --- [what was measured] (RQ2)}
  \begin{columns}[T]
    \column{0.55\textwidth}
      \placeholderfig{1}{0.45}{figure --- e.g.\ scaling plot}
    \column{0.45\textwidth}
      \scriptsize
      \begin{adjustbox}{max width=\linewidth}
      \begin{tabular}{@{}llrrr@{}}
        \toprule
        Case & $N$ & Size & Baseline & Proposed \\
        \midrule
        A1 &  64 & 7  &   7\,270 & 422 \\
        A2 & 256 & 9  & 119\,383 & 902 \\
        A3 & 1024 & 11 & ---     & 1\,614 \\
        B1 &  64 & 8  &  29\,655 & 153 \\
        B2 & 256 & 10 & ---      & 257 \\
        \bottomrule
      \end{tabular}
      \end{adjustbox}
      \vskip0.15cm
      \begin{itemize}
        \item baseline: its scaling
        \item proposed: \alert{its scaling} --- the headline saving
        \item the variant that does \emph{not} scale, and by how much
      \end{itemize}
  \end{columns}
  \scriptsize The condition under which the saving holds, and where it does not.
\end{frame}

% Result frame, layout B: two figures side by side, bullets below.
\begin{frame}{Result 2 --- [what was measured] (O4)}
  \begin{columns}[T]
    \column{0.5\textwidth}
      \placeholderfig{0.95}{0.45}{figure --- left}
    \column{0.5\textwidth}
      \placeholderfig{0.95}{0.45}{figure --- right}
  \end{columns}
  \small
  \begin{itemize}
    \item The reproduced known result, with the fitted exponent and the theorem it agrees with~\cite{example_article}.
    \item The extension, and what it costs.
    \item \alert{The one-sentence conclusion the committee should take from this frame.}
  \end{itemize}
\end{frame}

% Result frame, layout C: figure left, numeric table + bullets right.
\begin{frame}{Result 3 --- [what was measured] (O2)}
  \begin{columns}[T]
    \column{0.5\textwidth}
      \placeholderfig{1}{0.55}{figure}
    \column{0.5\textwidth}
      \scriptsize
      \begin{adjustbox}{max width=\linewidth}
      \begin{tabular}{@{}rrrrrl@{}}
        \toprule
        $N$ & $a$ & $r$ & $P(r)$ & cost & outcome \\
        \midrule
        15 & 7  & 4  & 0.51 &    522 & $3\times5$ \\
        21 & 2  & 6  & 0.31 & 10\,103 & $7\times3$ \\
        35 & 3  & 12 & 0.32 & 49\,739 & $7\times5$ \\
        \bottomrule
      \end{tabular}
      \end{adjustbox}
      \vskip0.15cm
      \scriptsize
      \begin{itemize}
        \item what the construction is, in one line
        \item what it achieves, quantitatively
        \item \alert{what it does not give}, stated as plainly as the positive result
      \end{itemize}
  \end{columns}
\end{frame}

% Result frame, layout D: figure left, table with a caption line above it.
\begin{frame}{Result 4 --- [what was measured] (O3)}
  \begin{columns}[T]
    \column{0.52\textwidth}
      \placeholderfig{1}{0.55}{figure}
    \column{0.48\textwidth}
      \scriptsize
      Setting; figure of merit defined in one line
      \begin{adjustbox}{max width=\linewidth}
      \begin{tabular}{@{}clrr@{}}
        \toprule
        $p$ & Method & Ratio & $P(\text{ok})$ \\
        \midrule
        1 & proposed (variant 1) & $0.852\pm0.022$ & 1.00 \\
          & proposed (variant 2) & $0.759\pm0.018$ & 1.00 \\
          & baseline             & $0.308\pm0.066$ & 0.38 \\
        4 & proposed (variant 1) & $0.910\pm0.020$ & 1.00 \\
          & proposed (variant 2) & $0.857\pm0.051$ & 1.00 \\
          & baseline             & $0.370\pm0.077$ & 0.47 \\
        \bottomrule
      \end{tabular}
      \end{adjustbox}
      \vskip0.15cm
      \scriptsize
      \begin{itemize}
        \item the construction and the identity or property it rests on
        \item the headline: \alert{$>2\times$} the baseline --- and \emph{why}
        \item honest reading: where the proposed method is weaker
      \end{itemize}
  \end{columns}
\end{frame}

% Result frame, layout E: figure left, model table with the proposed row in bold.
\begin{frame}{Result 5 --- [first learning / application result] (O5)}
  \begin{columns}[T]
    \column{0.5\textwidth}
      \placeholderfig{1}{0.55}{figure --- e.g.\ training curves}
    \column{0.5\textwidth}
      \tiny Task; metric over $k$ repetitions
      \setlength{\tabcolsep}{3pt}
      \begin{adjustbox}{max width=\linewidth}
      \begin{tabular}{@{}lrrr@{}}
        \toprule
        Model & Size & Params & Acc.\ (\%) \\
        \midrule
        reference model A               & 8 & 16 & $99.0\pm3.0$ \\
        \textbf{proposed (best variant)} & \textbf{5} & 32 & $\mathbf{99.0\pm2.0}$ \\
        proposed (variant 2)            & 5 & 32 & $92.7\pm4.4$ \\
        proposed (variant 3)            & 5 & 10 & $92.5\pm3.0$ \\
        reference model B               & 8 & 16 & $71.3\pm11.2$ \\
        \bottomrule
      \end{tabular}
      \end{adjustbox}
      \vskip0.15cm
      \scriptsize
      \begin{itemize}
        \item the reference reproduced, and what its result depends on
        \item the proposed model matches it \alert{at lower size} --- and why
        \item a variant that does worse, and the interpretation
        \item the cost caveat that keeps this a simulation result
      \end{itemize}
  \end{columns}
\end{frame}

\begin{frame}{Result 6 --- [second learning / application result] (O5)}
  \begin{columns}[T]
    \column{0.55\textwidth}
      \placeholderfig{1}{0.55}{figure --- e.g.\ accuracy by dataset}
    \column{0.45\textwidth}
      \scriptsize Secondary statistic, defined in one line
      \begin{adjustbox}{max width=\linewidth}
      \begin{tabular}{@{}lrrr@{}}
        \toprule
        Dataset & Proposed & Base.\ 1 & Base.\ 2 \\
        \midrule
        dataset 1 & 0.79 & 0.79 & 0.35 \\
        dataset 2 & 0.81 & 0.89 & 0.28 \\
        dataset 3 & 0.69 & 0.68 & \alert{0.10} \\
        dataset 4 & 0.60 & 0.60 & \alert{0.10} \\
        \bottomrule
      \end{tabular}
      \end{adjustbox}
      \vskip0.15cm
      \scriptsize
      \begin{itemize}
        \item a sanity check that passed
        \item the headline: \alert{best or joint-best among the like-for-like methods}
        \item the known failure mode~\cite{example_article} that the proposed method avoids, and why
        \item no advantage over the tuned classical baselines: say so
      \end{itemize}
  \end{columns}
\end{frame}

\begin{frame}{Result 7 --- [third learning / application result] (O5)}
  \begin{columns}[T]
    \column{0.5\textwidth}
      \placeholderfig{1}{0.55}{figure}
    \column{0.5\textwidth}
      \tiny Metric (\%), $k\times m$ folds
      \setlength{\tabcolsep}{2.5pt}
      \begin{adjustbox}{max width=\linewidth}
      \begin{tabular}{@{}lrrrr@{}}
        \toprule
        Method & Set 1 & Set 2 & Set 3 & Set 4 \\
        \midrule
        proposed        & $84.2\pm8.8$ & $57.0\pm7.0$ & $74.6\pm3.7$ & $70.6\pm3.5$ \\
        classical control & $87.2\pm8.3$ & $58.5\pm9.9$ & $74.8\pm4.3$ & $70.6\pm4.9$ \\
        standard baseline & $82.0\pm10.3$ & $60.3\pm6.4$ & $75.0\pm3.8$ & $71.0\pm3.9$ \\
        proposed, trained & $88.3\pm7.0$ & $57.6\pm6.6$ & --- & --- \\
        proposed, fixed   & $88.1\pm8.2$ & $57.5\pm4.9$ & --- & --- \\
        \bottomrule
      \end{tabular}
      \end{adjustbox}
      \vskip0.15cm
      \scriptsize
      \begin{itemize}
        \item what the proposed method reaches, against the standard baseline
        \item it does \alert{not} beat the classical control: say where and by how much
        \item an identified confound, and what must be repeated
      \end{itemize}
  \end{columns}
\end{frame}

\subsection[RQs Answered / Objectives Met]{How Research Questions are Answered and Research Objectives are Met}
% One row per case; the three RQ columns; \rqplus / \rqminus mark the
% direction of each finding.
\begin{frame}{How the research questions are answered (1): [first group of cases]}
  \scriptsize
  \renewcommand{\arraystretch}{1.25}
  \begin{tabularx}{\textwidth}{@{}>{\raggedright\arraybackslash\bfseries}p{0.13\textwidth} Y Y Y@{}}
    \toprule
    Case (obj.) & RQ1 --- formulation & RQ2 --- cost & RQ3 --- against a tuned baseline \\
    \midrule
    {}[Case A] (O1) & the formulation in one line & its cost & \rqplus\ the positive finding; \rqminus\ the limit \\
    {}[Case B] (O2) & the formulation in one line & its cost & the quantitative outcome; \rqminus\ what it does not give \\
    {}[Case C] (O3) & the formulation in one line & its cost & \rqplus\ the positive finding; \rqminus\ where it is weaker \\
    {}[Case D] (O4) & the formulation in one line & its cost & matches the known result; \rqminus\ the qualification \\
    \bottomrule
  \end{tabularx}
  \vskip0.2cm
  \textbf{O6} --- every negative outcome is reported at the same precision as the positive ones.
\end{frame}

\begin{frame}{How the research questions are answered (2): [second group of cases]}
  \scriptsize
  \renewcommand{\arraystretch}{1.25}
  \begin{tabularx}{\textwidth}{@{}>{\raggedright\arraybackslash\bfseries}p{0.13\textwidth} Y Y Y@{}}
    \toprule
    Case (obj.) & RQ1 --- formulation & RQ2 --- cost & RQ3 --- against a tuned baseline \\
    \midrule
    {}[Case E] (O5) & the formulation in one line & size vs.\ the reference; the depth caveat & \rqplus\ matched the reference; \rqminus\ a variant fits worse; \rqminus\ not a device proposal \\
    {}[Case F] (O5) & the formulation in one line & its cost & \rqplus\ best or joint-best like-for-like; \rqminus\ within noise of the classical baselines \\
    {}[Case G] (O5) & the formulation in one line & its cost & \rqplus\ matches the standard baseline; \rqminus\ never beats the classical control \\
    \bottomrule
  \end{tabularx}
  \vskip0.2cm
  \textbf{Regime identified (O6):} the one-sentence statement of where the approach pays and where it does not.
\end{frame}

\begin{frame}{Discussion: what [the approach] buys, and what it does not}
  \footnotesize
  \begin{columns}[T]
    \column{0.5\textwidth}
      \begin{block}{[The approach] buys}
        \begin{enumerate}
          \item \textbf{[Gain one]} --- the number that shows it.
          \item \textbf{[Gain two]} --- the number that shows it.
          \item \textbf{[Gain three]} --- the number that shows it.
          \item \textbf{[Gain four]}.
        \end{enumerate}
      \end{block}
    \column{0.5\textwidth}
      \begin{alertblock}{[The approach] does not buy}
        \begin{enumerate}
          \item {}[Non-gain one], stated plainly.
          \item {}[Non-gain two], with the size of the effect against the noise.
          \item {}[Non-gain three].
          \item {}[Non-gain four].
        \end{enumerate}
      \end{alertblock}
  \end{columns}
  \vskip0.2cm
  \begin{center}
    \textbf{Pattern:} the one-sentence rule that separates the two columns.
  \end{center}
\end{frame}

\begin{frame}{Threats to validity}
  \small
  \begin{itemize}
    \item \textbf{Scale.} The sizes reached, and what the fitted quantities do and do not establish.
    \item \textbf{Cost proxies.} What the cost measure ignores; which comparisons are therefore claimed.
    \item \textbf{Optimisation.} The optimiser and budget; what was \emph{not} measured --- name the largest single gap in the evidence.
    \item \textbf{Data.} Provenance of the datasets; which baselines reproduce published figures; what the synthetic tasks are and are not evidence for.
    \item \textbf{Scope.} The restrictions per case, in one line each.
  \end{itemize}
\end{frame}

% ================================================================
\section{Conclusion}
\begin{frame}{Conclusion}
  \small
  \begin{itemize}
    \item What was treated as \textbf{the object of study}, and across how many cases, on what terms.
    \item \textbf{The central lesson}, in one sentence.
    \item \textbf{Where it pays:} the list, in one line.
    \item \textbf{Where it does not:} the list, in one line.
    \item The scope of the claims; \textbf{what is not claimed}.
  \end{itemize}
\end{frame}

\begin{frame}{Limitations and future work}
  \footnotesize
  \begin{block}{Limitations}
    A single paragraph listing the restrictions of the study, separated by semicolons, with the most consequential one marked in italics.
  \end{block}
  \begin{block}{Future work (in order of how well-posed it is)}
    \begin{enumerate}
      \item The best-defined next step, with what it would cost.
      \item The measurement that was missing, with the references that define it~\cite{example_article, example_inproceedings}.
      \item The natural extension of the strongest result.
      \item The open question that decides whether the programme continues.
    \end{enumerate}
  \end{block}
\end{frame}

% ================================================================
\section{Misc Link (Git/Publication)}
\begin{frame}{Code, manuscript and publications}
  \begin{itemize}
    \item \textbf{Simulation code} (what it contains):\\
      \texttt{[repository URL]}
    \item \textbf{Manuscript source} (LaTeX, bibliography, generated tables):\\
      \texttt{[repository URL]}
    \item \textbf{Conference poster / talk} --- [venue, year]:\\
      \texttt{[title / link]}
    \item \textbf{Publication / preprint:}\\
      \texttt{[none yet / arXiv link]}
  \end{itemize}
  \vskip0.4cm
  \small One sentence on how the figures and tables in the talk were produced and how the random components were fixed.
\end{frame}

% ================================================================
\ifhasbib   % the whole reference list is skipped when there is no references.bib
\section*{References}
\begin{frame}[allowframebreaks]{References}
  \printbibliography[heading=none]
\end{frame}
\fi

\begin{frame}[plain]
  \begin{center}
    \vfill
    {\usebeamercolor[fg]{structure}\Huge Thank you}\\[0.5cm]
    {\large Questions}
    \vfill
  \end{center}
\end{frame}

\end{document}
